\documentclass{article}
\usepackage[margin=2cm]{geometry}
\usepackage{amsmath}
\usepackage{graphicx}
\usepackage{booktabs}
\usepackage[utf8]{inputenc}
\usepackage{ragged2e}
\setlength{\emergencystretch}{3em}
\begin{document}
\sloppy

\noindent\hfill\includegraphics[height=1.8em]{assets/figure_header_logo.png}\par\noindent\hrule\vspace{4pt}

is proportional to


\begin{equation}
\begin{array} { r } { \mathfrak { D } _ { 3 _ { D } } ^ { \Im _ { B }, ( 1, 0 ; A B ^ { - 1 } ) } T _ { \Im _ { D } } ( z ), } \\ { \mathrm { w h e r c ~ } \mathfrak { D } _ { \Im _ { D } } ^ { \Im _ { B }, ( 1, 0 ; A B ^ { - 1 } ) } \mathrm { ~ i s ~ g i v e n ~ b y ~ } ( ). } \end{array}
\end{equation}


\subsection*{1 Koornwinder limit of the trinion.}


We compute here the index of the three punctured sphere in the Koornwinder limit. The only subtle issue is the evaluation of the integrals in the
function  (??). The integral can be interpreted as index of two coupled$H$
SU(2) gauge theories with five flavors. The integrand does not have a good.
limit however by manipulating it using Seiberg dualities we can show that
the limit is well defined and evaluate all the integrals. We will discuss the.
evaluation of this function here. To evaluate  we first perform Seiberg du-$H$
ality on  splitting to five flavors in particular way, leading to SU(3) theory$u . { 2 }$
with five flavors


\begin{equation}
\begin{array} { r l } & { \quad \big ( ( q p ) ^ { \frac 1 4 } A ^ { - \frac 1 2 } t b z _ { 1 } ^ { \pm 1 }, ( q p ) ^ { \frac 1 4 } A ^ { \frac 1 2 } a ^ { - 1 } t v _ { 1 } ^ { \pm 1 }, ( q p ) ^ { \frac 1 4 } A ^ { - \frac 1 2 } b ^ { - 1 } z _ { 2 } ^ { - 1 } \big ), } \\ & { \quad \qquad \big ( ( q p ) ^ { \frac 1 2 } \frac 1 t ^ { 2 } w _ { 1 } ^ { \pm 1 },}\end{array}
\end{equation}


which leads to no$w . { 2 }$flavors having negative powers of  . We have the mesons$\pmb { p }$


\begin{equation}
\begin{array} { c } { ( q p ) ^ { \frac { 3 } { 4 } } A ^ { - \frac { i } { 2 } } t ^ { - 1 } b z _ { 1 1 } ^ { \pm 1 } w _ { 1 1 } ^ { \pm 1 }, ~ ( q p ) ^ { \frac { i } { 2 } } A ^ { \frac { 1 } { 2 } } a ^ { - 1 } t ^ { - 1 } w _ { 1 1 } ^ { \pm 1 }, ~ ( q p ) ^ { \frac { i } { 2 } } t ^ { - 2 } A ^ { - \frac { 1 } { 2 } } b ^ { - 1}}\end{array}
\end{equation}


and the new charged fields are.


sU(3)w (pq)A3b-3t3a-3z2\$z1, (pq)A-3b3t3a3z\$v1, (qp)A3b3t\$a-3z3
sU(3)w:(qp)t3Ata3b-3z\{w1, (qp)\$A-3b-3t-3a-3zvt1, (qp)A6b3t-sa3z2


and all of these fields have a good limit when goes to zero and the fugacities$p$
are scaled. We note that some of the mesons charged under  form mass$u . { 1 }$
terms with some of the quarks and after the first Seiberg duality we have
four flavors of$w . { 1 }$


\begin{equation}
\begin{array} { r l } & { \big \langle ( q p ) ^ { \frac 1 4 } A ^ { \frac 1 2 } b z _ { 2 }, ( q p ) ^ { \frac 1 1 2 } A ^ { \frac 1 4 } t ^ { \frac 2 3 } a ^ { - \frac 1 3 } z _ { 2 } ^ { \frac 1 3 } \big ( w _ { 2 } ^ { j } \big ) ^ { - 1 } \big \rangle, } \\ & { \big ( ( q p ) ^ { \frac 1 4 } A ^ { \frac 1 2 } b z _ { 2 } ^ { - 1 }, ( q p ) ^ { \frac 1 4 } A ^ { - \frac 1}}\end{array}
\end{equation}

\end{document}
